\hypertarget{link_to_nint}{}{\Large \bf nint}

{\it Integer nearest to the input shifted by a constant $c$}.


\begin{center}
\includegraphics{codegen/nint.pdf}
\end{center}
Syntax : \hspace*{10ex}\parbox[t]{\textwidth}{\sf 
\& nint \\
name of variable $x_i$\\
name of variable $x_j$\\
data name, parameter name or math expression for $c$
}

~ \hrule ~

Internal states : none

Discrete variables : $z \in {\cal I}$

Equations :
\begin{equation*}
0 = x_j - z
\end{equation*}

Discrete transitions : 

\hspace*{20ex}\begin{minipage}[t]{15cm}
\begin{algorithmic} 
\IF {$\mbox{nint}(x_i+c) \ne z$} 
  \STATE $z \gets \mbox{nint}(x_i+c)$
\ENDIF
\end{algorithmic}
\end{minipage}

~

where the $\mbox{nint}$ function returns the nearest integer.

Particular cases:
\begin{itemize}
\item with $c=0$, the output is the integer nearest to $x_i$;
\item with $c=-0.5$, the output is the ``floor'' integer of $x_i$, i.e. the largest integer smaller or equal to $x_i$;
\item with $c=0.5$, the output is the ``ceiling'' integer of $x_i$, i.e. the smallest integer larger or equal to $x_i$.
\end{itemize}

\vspace*{2ex}
Initialization of discrete variables :
\[ z \gets \mbox{nint}(x_i+c) \]